% Source-derived chapter 03: Lagrangian fibrations
% Original source: riemann-roch-polynomials-hyperkahler-fourfolds
\chapter{Lagrangian fibrations}
\label{riemann-roch-polynomials-hyperkahler-fourfolds-chapter-03}
\source{riemann-roch-polynomials-hyperkahler-fourfolds}

\begin{remark}[Source section]
\label{riemann-roch-polynomials-hyperkahler-fourfolds-chapter-03-source}
\source{riemann-roch-polynomials-hyperkahler-fourfolds}
\source{riemann-roch-polynomials-hyperkahler-fourfolds}
This chapter reproduces the corresponding section of the retrieved
LaTeX source, including its definitions, statements, examples, and
proofs. It has not yet been translated into Lean.
\notready
\end{remark}

% The source section title was: Lagrangian fibrations
\label{sec:LagrangianFibration}
\source{riemann-roch-polynomials-hyperkahler-fourfolds}


Let $X$ be again a \hk\ manifold of dimension $2n$.\
We assume in this section that there is a  fibration $f\colon X\to \P^n$.\
Set $L\coloneqq f^*\cO_{\P^n}(1)$, with first Chern class $\lll\in H^2(X,\Z)$; it satisfies \mbox{$q_X(\lll)=0$.}\
Let $\mm\in H^2(X,\Z)$ be another class, not necessarily of type $(1,1)$; we assume \mbox{$q_X(  \lll, \mm)>0$.}

Any smooth fiber $X_b\coloneqq f^{-1}(b)$ is a Lagrangian complex torus  of dimension~$n$ (\cite[Theorem~1]{mat}, \cite[Theorem 1]{ac}).\
By \cite[Lemma~2.2]{matlag}, the hyperplane $\lll^\bot$ is contained in the kernel of the  restriction map $r_b\colon H^2(X,\C)\to H^2(X_b,\C)$.\
Since the restriction of a K\"ahler class on $X$ is a K\"ahler class on $X_b$, the map $r_b$ has rank exactly~$1$ and the rational class $r_b(\mm)$ is  a positive multiple of a K\"ahler class, hence is an ample class on $X_b$.\
In particular, $X_b$ is an abelian variety (\cite[Proposition~4]{ac}) and the ``degree'' of the polarization $r_b(\mm)=\mm\vert_{X_b}$ is the positive integer
\begin{equation*}
  \frac{1}{n!}\int_{X_b}(\mm\vert_{X_b})^{n}= \frac{1}{n!}\int_X \lll^{n}\mm^{n}=a
\end{equation*}
already considered in~\eqref{defa2}.



There are restrictions on the values that $a$ can take (see Theorem~\ref{th43}   for restrictions on small values of $a$ when $n=2$).\
For the moment, we prove that the existence of the Lagrangian fibration $f$ imposes strong conditions on~$X$ when~$a=1$, that is, when there is a class on $X$ that induces a principal polarization on the smooth fibers of $f$.

\subsection{The Huybrechts--Riemann--Roch polynomial}\label{subsec:RRpolynomialLagrangian}
\source{riemann-roch-polynomials-hyperkahler-fourfolds}

Keeping the notation and the hypotheses as above, we show that in the case $a=1$, the Huybrechts--Riemann--Roch polynomial of $X$ is completely determined.\
The main idea of the proof is taken from~\cite{ort}: the hypothesis~$a=1$ is essential because it implies that a certain locally free sheaf of rank $a$ on~$\P^n$ can be written as $\cO_{\P^n}(d)$ for some integer $d$.

\begin{theo}\label{th21}
\source{riemann-roch-polynomials-hyperkahler-fourfolds}
Let $X$ be a \hk\ manifold of dimension~$2n$ with a Lagrangian fibration $f\colon X\to \P^n$.\
Set $\lll\coloneqq c_1(f^*\cO_{\P^n}(1))\in H^2(X,\Z)$ and assume that there exists $\mm\in H^2(X,\Z)$ such that $\int_X \lll^{n}\mm^{n}=n!$.\
Then, $q_X(\lll,\mm)=\pm 1$, the quadratic form~$q_X$ is even, $c_X=(2n-1)!!$,
\begin{equation*}
P_{RR,X}(T)= \binom{\frac{T}{2}+n+1}{n},
\end{equation*}
and the sublattice $\Z\lll\oplus \Z\mm$ of $(H^2(X,\Z),q_X)$ is isomorphic to a hyperbolic plane.
\end{theo}

For all known \hk\ manifolds $X$, the lattice $(H^2(X,\Z),q_X)$ contains a hyperbolic plane.

\begin{proof}
Changing $\mm$ into $-\mm$ is necessary, we may assume $q_X(\lll,\mm)>0$.\
Consider the universal family $(\cX,\cL)\to {\mathfrak M}_\lll$ over one component of the (non-Hausdorff) moduli space of marked hyper-K\"ahler manifolds with a fixed $(1,1)$-class $\lll$.\
The period map $\cP\colon{\mathfrak M}_\lll\to \PP(\lll^\perp\otimes\C)$ is injective over very general points and the points in an arbitrary fiber correspond to the chambers of the decomposition of the positive cone.\
For a distinguished point $0\in {\mathfrak M}_\lll$ we have $(\cX_0,\cL_0)\cong(X,L )$, where $L\coloneqq f^*\cO_{\P^n}(1)$.\
But there also exists a fiber $X^\prime\coloneqq\cX_{0^\prime}$, with $0^\prime\in{\mathfrak M}_\lll$, whose rational N\'eron--Severi group~$\NS(X^\prime)_\Q$ is generated by the classes $\lll$ and $\mm$.\
In fact, since the lattice generated by $\lll$ and~$\mm$ is saturated (Remark~\ref{prim}), the integral N\'eron--Severi group  is generated by  $\lll$ and $\mm$.\
Furthermore, replacing $0^\prime$ by another point in the same fiber over $\cP(0^\prime)$, we may assume that $L^\prime\coloneqq\cL_{0^\prime}$ is nef.\
Since $q_X(k\lll+\mm)=2kq_X( \lll,\mm)+q_X( \mm)>0$ for $k\gg 0$, the manifold  $X'$ is projective by \cite{HuyHK}.

We apply~\cite[Theorem 1.2, Claim 3.1 and Claim 3.2]{mat:isotropic}: there exists a Lagrangian fibration $f^\prime\colon X^\prime\to \P^n$ such that $f^{\prime\,*}\cO_{\P^n}(1)\cong L^\prime$.\
Upon replacing $(X,L)$ by $(X^\prime,L^\prime)$, we may, since $q_X$, $c_X$, and $P_{RR,X}(T)$ are invariant by deformation, assume that $X$ carries a line bundle~$M$ with first Chern class $\mm$.

Since $L$ is nef, $q_X(\lll,\mm)>0$, and $\NS(X)$ is generated by $\lll$ and $\mm$, we can replace $M$ with $L^k\otimes M$, for $k\gg0$, and assume that $M$ is ample on $X$.\  By \cite[Theorem~10.32]{kol}, $R^if_*M$ vanishes for $i>0$.\ Because $f$ is flat,     the sheaf $\cM\coloneqq f_*M$ on~$\P^n$ is locally free; since
the restriction of $M$ to a smooth generic fiber of $f$ defines a principal polarization, the rank of $\cM$ is  $a=1$.\
By the projection formula and~\eqref{rr}, $\cM$ satisfies
\begin{equation}\label{chi}
\source{riemann-roch-polynomials-hyperkahler-fourfolds}
\forall k\in\Z\qquad\chi(\P^n, \cM(k))=\chi(X,  L^k\otimes M) =P_{RR,X}(2kq_X(\lll,\mm)+q_X(\mm)) .
\end{equation}
Following \cite[Section~3]{ort}, we write $\cM=\cO_{\P^n}(d)$ for some integer $d$ and, from~\eqref{chi}, we deduce
\[
\forall k\in\Z\qquad P_{RR,X}(2kq_X(\lll,\mm)+q_X(\mm))=\chi(\P^n,\cO_{\P^n}(d+k))=\binom{d+k+n}{n},
\]
so that
\begin{equation}\label{bin}
\source{riemann-roch-polynomials-hyperkahler-fourfolds}
P_{RR,X}(T)= \binom{d+\frac{T-q_X(\mm)}{2 q_X(\lll,\mm)}+n}{n}.
\end{equation}
Set $\gamma\coloneqq \frac{ q_X(\mm)}{2 q_X(\lll,\mm)} $.\ Since $P_{RR,X}(0)=n+1$, either $d-\gamma=1$, or $n$ is even and $d-\gamma=-n-2$.\footnote{The equation $\binom{x+n}{n}= n+1$ is equivalent to the monic equation $\prod_{i=1}^n(x+i)= (n+1)! $, so any rational solution is in fact integral, and the only integral solutions are $x=1$ and, when $n$ is even, $x=-n-2 $.}\
Since the coefficients of $P_{RR,X}$ are positive and the coefficient of $T^{n-1}$ in~\eqref{bin} is a positive multiple of $ n+2(d-\gamma) +1$, the latter case is ruled out.\
So $\gamma$ is an integer and $d=\gamma+1$.\ Replacing $M$ by $L^{-\gamma}\otimes M$, we may  assume $\gamma=0$ and $d=1$, hence $q_X(\mm)=0$, so that~\eqref{bin} becomes
\begin{equation}\label{bin2}
\source{riemann-roch-polynomials-hyperkahler-fourfolds}
P_{RR,X}(T)= \binom{\frac{T}{2 q_X(\lll,\mm)}+n+1}{n}.
\end{equation}
By  Lemma~\ref{star} below (applied with $c=n!$, $c'=1$, and $q=2 q_X(\lll,\mm)$), we get    $q_X(\lll,\mm)=1$ and the quadratic form~$q_X$ is even.\
The value of $c_X$ is then derived from~\eqref{pol} and the polynomial $P_{RR,X}(T)$ from~\eqref{bin2}.
\end{proof}

It remains to prove the arithmetical result   used at the end of the proof above.\ We prove a bit more  than what we actually used above but we will need this stronger statement for the proof of Theorem~\ref{th43}.

\begin{lemm}\label{star}
\source{riemann-roch-polynomials-hyperkahler-fourfolds}
Let $X$ be a \hk\ manifold of dimension $2n$.\ Assume that there are positive integers $c$, $c'$, and $q$  such that $c'$ is divisible by no nontrivial $n$th powers and the polynomial $P(T)\coloneqq \frac{c}{c'}P_{RR,X}(qT)$ is  monic  with integral coefficients.\ Then, either $q=1$, or $q=2$  and the quadratic form $q_X$ is even.
\end{lemm}

\begin{proof}
Let $\alpha\in H^2(X,\Z)$ and write $\frac{ q_X(\alpha)}{q}=:\frac{r}{s}$, where $r$ and $s$ are relatively prime integers.\ One has
\begin{equation*}
    P_{RR,X}( q_X(\alpha))=\frac{c'}c\, P\Bigl(\frac{q_X(\alpha)}{q}\Bigr)
    =\frac{c'}c\, P\Bigl(\frac{r}{s}\Bigr)
\end{equation*}
and this is an integer by Lemma~\ref{lem11}.\ Since $P(T)$ is monic  with integral coefficients, $s^nP(\frac{r}{s})$ is an integer congruent to $r^n$ modulo $s$, hence prime to $s^n$.\ But $c's^nP(\frac{r}{s})= cs^nP_{RR,X}( q_X(\alpha))$ is divisible by $s^n$, hence so is $c'$.\ Our hypothesis implies $s=\pm1$, which proves that  $q$ divides all values that $q_X$ takes on $H^2(X,\Z)$.\
Since the  integral bilinear form associated with $q_X$ is  not  divisible, either $q=1$, or $q=2$  and the quadratic form $q_X$ is even.
\end{proof}


\subsection{The exceptional divisor}\label{subsec:ExceptionalDivisor}
\source{riemann-roch-polynomials-hyperkahler-fourfolds}

We keep our \hk\ manifold $X$ of dimension $2n$ with a Lagrangian fibration $f\colon X\to \P^n$ and we set as above $L\coloneqq f^*\cO_{\P^n}(1)$, with first Chern class $\lll\in H^2(X,\Z)$.\
We assume further that there exists a line bundle $M$ on $X$ whose class~$\mm$ satisfies $\int_X \lll^{n}\mm^{n}=n!$ and $q_X(\lll,\mm)>0$.\

As in the proof of Theorem~\ref{th21}, $X$ is projective and we may assume  $q_X(\lll,\mm)=1$ and $q_X(\mm)=0$.\
Then $f_*M=\cO_{\P^n}(1)$ and $f_*(L^{-1}\otimes M )=\cO_{\P^n}$, hence the linear system  $| L^{-1}\otimes M |$ contains a single (effective) divisor~$E$  which still induces a principal polarization on the smooth fibers of $f$.\
It satisfies $q_X([E])=-2$ and $q_X([E],\mm)=-1$; in particular, $M$ is not nef.

\begin{lemm}\label{lem24}
\source{riemann-roch-polynomials-hyperkahler-fourfolds}
Assume that $\NS(X)=\Z \lll \oplus \Z \mm$.
Then the divisor $E\in | L^{-1}\otimes M |$ is irreducible and reduced.
\end{lemm}

\begin{proof}
For any integer $k\ge 2$, the linear system $| L^{-k}\otimes M |$ is empty, because any divisor in that linear system would have negative $q_X$-intersection with $E$, and would therefore contain~$E$, but the linear system $|L^{-k+1}|$ is empty.\
It follows that $E$ has no vertical components.\ The fact that $E$ is   irreducible and reduced then follows from the relation $q_X([E],\lll)=1$.
\end{proof}

Since $q_X( [E])<0$, the prime divisor $E$ is therefore exceptional in the sense of \cite[d\'ef.~3.10]{bou} and
\cite[Definition~3.2]{mar} hence it spans an extremal ray in the effective cone $\Eff(X)$.\
Moreover, by~\cite[prop.~1.4 and rem.~4.3]{dru}, there is a birational isomorphism $\phi\colon X\isomdra X' $ (with~$X'$ smooth \hk) and a projective divisorial contraction $c\colon X'\to Y$ with exceptional divisor $\phi(E)$; moreover, the general fibers of $c\vert_{\phi(E)}$ are either smooth rational curves, or unions of two smooth rational curves meeting transversely at one point.\
The divisor $E$ is uniruled, its class $-\lll+\mm$ is primitive (because $q_X(\lll,[E])=1$), the reflection
\[
\alpha \longmapsto  \alpha+ q_X(\alpha,-\lll+\mm) (-\lll+\mm)
\]
is integral and a monodromy operator that permutes $\lll$ and $\mm$.\
Finally,  the class in $H_2(X',\Z)\isom H^2(X',\Z)^\vee\isom H^2(X,\Z)^\vee$ of a general (curve) fiber of $\phi(E)\to c(\phi(E))$ is given by the linear form $q_X(-\lll+\mm,\bullet)$; moreover, since $q_X(-\lll+\mm,\lll)=1$, this general fiber cannot be the union of two homologous curves, hence it is a smooth rational curve (\cite[Corollary~3.6]{mar}).\


\subsection{Examples}\label{subsec:SYZKnownExs}
\source{riemann-roch-polynomials-hyperkahler-fourfolds}
 The next two examples show that Theorem~\ref{th21} applies to  \hk\ manifolds of K3$^{[n]}$ deformation type (Example~\ref{ex:K3n}) or of $\mathrm{OG10}$ deformation type (Example~\ref{ex:OG10}).

\begin{exam}
\source{riemann-roch-polynomials-hyperkahler-fourfolds}
\notready\label{ex:K3n}
\source{riemann-roch-polynomials-hyperkahler-fourfolds}
Let $S$ be a K3 surface with a primitive polarization $\hh_S$ of degree $\hh_S^2=2d$ and set $n=d+1$ (this is the genus of any curve in the linear system $ |\hh_S|$).\
Assume that the pair $(S,\hh_S)$ is very general, so that $\NS(S)=\Z \hh_S$.\
The smooth projective moduli space  $\M_0(S)\coloneqq M_{S}(0,\hh_S,0)$ parametrizes pairs consisting of a curve $C\in|\hh_S|$ and a torsion-free, rank-1 coherent sheaf on~$C$ of degree $n-1$, considered as a (torsion) sheaf on $X$.\
It is a \hk\ variety of K3$^{[n]}$ deformation type.\
There is a Lagrangian fibration
$
f\colon \M_0(S)\lra |\hh_S|=\P^{n}
$
that takes a sheaf to its support.\
The fiber of a smooth $C\in  |\hh_S|$ is the Jacobian $J^{n-1}(C)$.

The lattice $\NS(\M_0(S))$   is spanned by two isotropic vectors
$\lll\coloneqq c_1(f^*\cO_{\P^{n}}(1))$ and $  \mm$
 which satisfy $ q_X(\lll, \mm)=1$.\  Since the Fujiki constant is $ (2n-1)!!$, formula~\eqref{defia} gives $a = q_X(\lll, \mm)^n=1$.\  The fibers $J^{n-1}(C)$ have canonical theta divisors that fit together to define an effective divisor in $\M_0(S)$ which is the exceptional divisor $E$ from Section~\ref{subsec:ExceptionalDivisor}.
\end{exam}

\begin{exam}[R\'ios Ortiz]
\source{riemann-roch-polynomials-hyperkahler-fourfolds}
\notready\label{ex:OG10}
\source{riemann-roch-polynomials-hyperkahler-fourfolds}
According to~\cite{lsv}, there is a \hk\ manifold~$X$ of OG10 deformation type with a Lagrangian fibration \mbox{$f\colon X\to \P^5$} and an $f$-ample effective divisor~$\Theta$ on~$X$ that restricts to a principal polarization on the smooth fibers of $f$ (\cite[Proposition~5.3]{lsv}).\
Set $L\coloneqq f^*\cO_{\P^n}(1)$ and $M\coloneqq\cO_X(\Theta)$.\
Theorem~\ref{th21} shows that the Fujiki constant $c_X$ is equal to $9!!=945$ (it was originally computed in~\cite{rap}) and that
$
P_{RR,X}(T)= \binom{ \frac{T }{2  }+6}{5}$ (\cite{ort}).
\end{exam}

The next two examples deal with the other   known types of  \hk\ manifolds.

\begin{exam}
\source{riemann-roch-polynomials-hyperkahler-fourfolds}
\notready\label{ex:Kumn}
\source{riemann-roch-polynomials-hyperkahler-fourfolds}
 Let $A$ be an abelian surface with a polarization~$\hh_A$ of type $(1,d)$, with $ d\ge 3$.\
Assume that the pair $(A,\hh_A)$ is very general, so that $\NS(A)=\Z \hh_A$.\
The smooth projective moduli space  $\M_0(A)\coloneqq M_{A}(0,\hh_A,0)$ parametrizes pairs consisting of a curve $C\subset A$ with class~$\hh_A$ and a torsion-free, rank-1 coherent sheaf on~$C$ of degree $d$.\
The fibers of its (surjective) Albanese map $\M_0(A)\to A\times \widehat A$   are all isomorphic to the same \hk\ manifold $\K_0(A)$ of dimension $2n\coloneqq 2d-2$ which is of generalized Kummer deformation type (\cite[Theorem~0.2(1)]{yos:moduli}).\
There is a Lagrangian fibration
$
f\colon \K_0(A) \to \P^{n}
$
that takes a sheaf to its support.\
The fiber of a smooth $C\num \hh_A$ is the kernel of the Abel--Jacobi map $J^{d}(C)\to A$.\

By~\cite[Theorem 0.2(2)]{yos:moduli}, the lattice $\NS(\K_0(A))$ is spanned by the two isotropic vectors $\lll\coloneqq c_1(f^*\cO_{\P^{n}}(1))$ and $  \mm=c_1(M)$ which satisfy $q_X(\lll, \mm)=1$.\
Since the Fujiki constant is $(n+1) (2n-1)!!$, formula~\eqref{defia} gives $a=n+1$.
\end{exam}

\begin{exam}
\source{riemann-roch-polynomials-hyperkahler-fourfolds}
\notready
Examples of \hk\ manifolds of  OG6 deformation type with a Lagrangian fibration are described in \cite{rap}.
%: let $(A,\theta_A)$ be a principally abelian surface such that $\NS(A)=\Z\theta_A$ and let  $\M(A):=M_{A}(0,2\theta_A,-2)$ be the (singular) projective moduli space   parametrizing pairs consisting   of a curve $C\subset A$ with class $2\theta_A$ and   a torsion-free, rank-1 coherent sheaf on~$C$ with Euler characteristic $-2$.\ The fiber $\K(A) $ over $(0,0) $ of its surjective Albanese map $\M(A)\to A\times \widehat A$ has a canonical desingularization $\widetilde\K(A) $ which is a  \hk\ manifold   of OG6 deformation type (\cite[Corollary~1.1.8 and Proposition~2.2.1]{rap}).\ There is a canonical Lagrangian fibration $f\colon \widetilde\K(A)\to \K(A)\to |2\Theta_A|=\P^3$ that sends a sheaf to its support (\cite[Section~2.1]{rap}).\ The fiber of a smooth $C\num 2\ell_A $ is the kernel of the Abel--Jacobi map  $J^2(C)\to A$.\
\end{exam}

%%%%%%%%%%%%%%%%%%%%%
